\section{Motivation}
Research today is fragmented across disconnected tools by default, not by
choice. The survey behind the ScholarFlow SRS quantified the landscape:
34.4\% of researchers read directly in the browser with no organization,
31.3\% rely on local folders with no metadata, 31.3\% store files in cloud
drives that hold files but no paper context, 28.1\% use no tool at all, and
only 15.6\% use Zotero --- the legacy citation manager whose pre-AI interface
and basic group libraries disappoint modern expectations shaped by Notion,
Figma, and ChatGPT.

Three market failures explain the fragmentation. First, \emph{legacy tool
stagnation}: reference managers designed between 2006--2011 never absorbed
the transformer revolution, so papers are still searched by keywords instead
of meaning, summaries are still manual, and citations still lack context.
Second, \emph{workflow disintegration}: a researcher cobbles together Drive
for storage, Zotero for citations, Notion for notes, email for collaboration,
and ChatGPT for summaries via manual copy-paste --- five silos whose
context-switching overhead cuts research velocity by an estimated 30--50\%.
Third, \emph{the AI integration gap}: no platform treats AI as a first-class
citizen of the research workflow rather than a bolt-on chatbot.

The motivation for ScholarFlow is therefore threefold: modernize the research
ecosystem with one accessible, transparent, efficient platform; deliver
jurisdiction- and context-aware AI assistance (summaries, Q\&A, comparison,
literature synthesis) grounded in the user's own library instead of generic
web text; and do it at student-friendly pricing, since professional legacy
tiers (\$50--\$250/user/month) exclude exactly the students and small groups
who need leverage most. The timing is favorable --- LLM quality and cost now
permit everyday summarization and synthesis, pgvector brings semantic search
into ordinary Postgres, and modern web frameworks let small teams ship
real-time collaborative applications --- and the survey confirms the demand:
75\% need validation with zero respondents uninterested in AI features.
